\documentclass[tikz,border=4pt]{standalone}
\usepackage{tikz}
\usetikzlibrary{arrows.meta,positioning,shapes.geometric,backgrounds,fit,calc}

\begin{document}
\begin{tikzpicture}[
  font=\small,
  >=Stealth,
  layerbox/.style={
    rectangle, rounded corners=5pt,
    minimum width=10cm, text width=9.8cm,
    inner sep=6pt,
    draw=#1!70!black, fill=#1!10,
    align=center
  },
  agentbox/.style={
    rectangle, rounded corners=3pt,
    draw=blue!60!black, fill=blue!20,
    minimum width=1.45cm, minimum height=0.7cm,
    text centered, text width=1.4cm, font=\scriptsize
  },
  subbox/.style={
    rectangle, rounded corners=3pt,
    draw=teal!60!black, fill=teal!15,
    minimum width=2.8cm, minimum height=0.65cm,
    text centered, text width=2.7cm, font=\scriptsize
  },
  iobox/.style={
    rectangle, rounded corners=3pt,
    draw=gray!60, fill=gray!12,
    minimum width=3.4cm, minimum height=0.65cm,
    text centered, text width=3.3cm, font=\scriptsize
  },
  arr/.style={->, thick, gray!70},
]

%% --- Layer 1: User Interaction ---
\node[layerbox=gray, label={[font=\scriptsize\bfseries]left:}] (L1) at (0,0) {
  \textbf{User Interaction Layer}\\[3pt]
  \begin{tikzpicture}[baseline]
    \node[iobox] (inp) {\textbf{Input:} Backend / DSA\\Problem Statement};
    \node[iobox, right=1.5cm of inp] (btn) {\textit{Synthesize Logic}};
    \draw[arr] (inp) -- (btn);
  \end{tikzpicture}
};

%% --- Layer 2: Workflow Orchestration ---
\node[layerbox=orange, below=0.25cm of L1] (L2) {
  \textbf{Workflow Orchestration Layer}\\[3pt]
  \begin{tikzpicture}[baseline]
    \node[subbox, minimum width=8cm, text width=7.9cm]
      {\textbf{LangGraph Orchestrator}\\
       Manages Agent Execution, Conditional Transitions,\\
       Retry Mechanisms, State-Aware Flow};
  \end{tikzpicture}
};

%% --- Layer 3: Agent Layer ---
\node[layerbox=blue, below=0.25cm of L2] (L3) {
  \textbf{Agent Layer}\\[4pt]
  \begin{tikzpicture}[baseline]
    \node[agentbox]                        (A1) {Planning\\Agent};
    \node[agentbox, right=0.3cm of A1]     (A2) {Code\\Generation\\Agent};
    \node[agentbox, right=0.3cm of A2]     (A3) {Logical\\Review\\Agent};
    \node[agentbox, right=0.3cm of A3]     (A4) {Evaluation\\Agent};
    \node[agentbox, right=0.3cm of A4]     (A5) {Optimization\\Agent};
    \node[agentbox, right=0.3cm of A5]     (A6) {Explanation\\Agent};
    \foreach \a/\b in {A1/A2,A2/A3,A3/A4,A4/A5,A5/A6}
      \draw[->,thick,blue!50] (\a) -- (\b);
    %% feedback arrow
    \draw[->,thick,red!60] (A4.north) -- ++(0,0.3) -| node[above,font=\tiny,text=red!60]{Execution Feedback} (A2.north);
  \end{tikzpicture}
};

%% --- Layer 4: Execution and Evaluation ---
\node[layerbox=teal, below=0.25cm of L3] (L4) {
  \textbf{Execution and Evaluation Layer}\\[3pt]
  \begin{tikzpicture}[baseline]
    \node[subbox] (sb) {Code Interpreter\\/ Sandbox};
    \node[subbox, right=1.0cm of sb] (ts) {Test Suite\\Runner};
    \draw[arr] (sb) -- (ts);
  \end{tikzpicture}
};

%% --- Layer 5: Optional Knowledge Memory ---
\node[layerbox=purple, below=0.25cm of L4] (L5) {
  \textbf{Optional Knowledge Memory Layer}\\[3pt]
  \begin{tikzpicture}[baseline]
    \node[subbox, draw=purple!60!black, fill=purple!15, minimum width=5cm, text width=4.9cm]
      {ChromaDB Vector Store\\(Algorithmic Techniques \& Patterns)};
  \end{tikzpicture}
};

%% --- Output panel label ---
\node[iobox, right=0.5cm of L2, draw=green!60!black, fill=green!10,
      minimum height=4cm, text width=2.6cm, align=left, font=\scriptsize] (out) {
  \textbf{Output Dashboard}\\[4pt]
  Generated Code \& Plan\\[2pt]
  Execution Results\\(Pass/Fail, Tests)\\[2pt]
  System Metrics \&\\Confidence Score\\[2pt]
  Severity Analysis\\[2pt]
  Execution Trace Logs
};

%% Inter-layer arrows
\foreach \a/\b in {L1/L2, L2/L3, L3/L4, L4/L5}
  \draw[arr] (\a.south) -- (\b.north);

%% RAG query arrow from Agent layer to Memory
\draw[->,dashed,purple!60,thick] (L3.south west) -- ++(0,-0.25) -- ++(0,-0.25) |- (L5.west)
  node[near start, left, font=\tiny, text=purple]{Query/Retrieve};

\end{tikzpicture}
\end{document}
